%% kapa: exam=1 round=34 year=2023 period=2 form=A subject=appraisal_law subject_ko=감정평가관계법규 seq=13 number=13
%% source: Q-Net 감정평가사 기출문제 (cst00309) artlSeq=5211789 file=raw/1008/5211789/2023년도 제34회 감정평가사 1차시험 2교시.pdf
%% conversion: kordoc 4.24.51 -> tools/gen1.py (auto); review=proofread-sonnet; flags=none
\begin{question}{13}
국토의 계획 및 이용에 관한 법령상 국토교통부장관, 시·도지사, 시장 또는 군수나 도시·군계획시설사업의 시행자가 타인의 토지에 출입할 수 있는 경우를 모두 고른 것은?
\begin{framedtext}
ㄱ. 도시·군계획에 관한 기초조사 \hfill ㄴ. 개발밀도관리구역에 관한 기초조사 \par
ㄷ. 도시·군계획시설사업에 관한 측량 \hfill ㄹ. 지가의 동향에 관한 조사 \par
\end{framedtext}
\begin{choices}
  \item ㄱ
  \item ㄴ
  \item ㄷ, ㄹ
  \item ㄱ, ㄴ, ㄷ
  \item ㄱ, ㄴ, ㄷ, ㄹ
\end{choices}
\end{question}
\answer{5}
